\documentclass[11pt,a4paper]{article}
\usepackage[margin=1in]{geometry}
\usepackage{amsmath,amssymb,amsthm}
\usepackage{algorithm}
\usepackage{algpseudocode}
\usepackage{booktabs}
\usepackage{hyperref}

\theoremstyle{definition}
\newtheorem{definition}{\textbf{Definition}}[section]
\newtheorem{theorem}{\textbf{Theorem}}[section]
\newtheorem{lemma}[theorem]{\textbf{Lemma}}
\newtheorem{remark}{\textbf{Remark}}[section]

\title{\textbf{FP8 Mixed-Precision Training: E4M3/E5M2 Formats and Delayed Scaling}}
\author{\textbf{Team Scaibu} \\ \small Email: \texttt{jaydeep.9664920749@gmail.com} \\ \small \textbf{Architecture and Core Engine Team}}
\date{\textbf{October 2026}}

\begin{document}
\maketitle

\section{Definitions and Cognitive Mental Models}

\subsection{Formal Mathematical Definition}
\textbf{Definition (8-Bit Floating Point Formats and Delayed Dynamic Scaling):}
FP8 training (Micikevicius et al., 2022) leverages 8-bit floating-point arithmetic on modern tensor core architectures (NVIDIA Hopper H100/Blackwell B200) to double computational throughput and halve memory bandwidth relative to 16-bit half precision.

The Open Compute Project (OCP) standardizes two distinct 8-bit formats:
\begin{enumerate}
    \item \textbf{E4M3 (1 sign, 4 exponent, 3 mantissa bits)}: Higher precision, lower dynamic range ($[-448, 448]$). Primary format used for forward pass activations and weights.
    \item \textbf{E5M2 (1 sign, 5 exponent, 2 mantissa bits)}: Lower precision, identical exponent range to FP16 ($[-57344, 57344]$). Primary format used for backward pass gradients.
\end{enumerate}

Because FP8 formats possess narrow representable intervals, tensors $X \in \mathbb{R}^{d_1 \times d_2}$ are dynamically scaled by factor $S > 0$:
{\boldmath
\begin{equation}
X_{\text{FP8}} = \operatorname{clip}\left( \operatorname{round}(S \cdot X), \, -X_{\max}^{\text{FP8}}, \, X_{\max}^{\text{FP8}} \right)
\end{equation}
}
Under \textbf{Delayed Scaling}, calculating the exact maximum absolute value ($a_{\max} = \max |X|$) of the current tensor before every GEMM would introduce a synchronization barrier. Instead, the scaling factor for iteration $t$ is computed from a rolling history window of past $a_{\max}$ values:
{\boldmath
\begin{equation}
S_t = \frac{X_{\max}^{\text{FP8}}}{\max_{k \in [1, H]} a_{\max}^{(t-k)} \cdot 2^M}
\end{equation}
}
where $H$ is the history window length (e.g. $H=16$) and $M \ge 0$ is a headroom margin preventing unexpected overflow.

\subsection{Expanded Non-Technical Definition (Intuition for Anyone)}
\begin{enumerate}
    \item \textbf{Fitting a Giant Sound Wave on an 8-Bit Recorder}: An 8-bit number has only 256 possible values! If you record a whisper or an explosion without adjusting the microphone gain, the audio either has massive static or clips completely. FP8 scaling adjusts the volume gain knob before every layer so the signal uses all 256 levels perfectly.
    \item \textbf{Two Specialized Tools for Two Jobs}:
    \begin{itemize}
        \item \textbf{E4M3}: High resolution for thoughts (weights and forward thoughts) where precision matters.
        \item \textbf{E5M2}: Wide dynamic range for corrections (backward gradients) that can vary by orders of magnitude.
    \end{itemize}
    \item \textbf{Zero Synchronization Latency (Delayed Scaling)}: Instead of pausing all 1,000 GPUs to find the highest number in the current tensor, GPUs use the highest number from the last few steps, keeping hardware pipelines moving at maximum speed.
\end{enumerate}

\subsection{Child-Friendly Mental Model (Physical Metaphor)}
Imagine packing books into a tiny backpack with only 8 pockets:
\begin{itemize}
    \item \textbf{Without Scaling}: Trying to fit encyclopedia volumes into mini pockets. Everything rips!
    \item \textbf{FP8 Scaling}: Taking high-resolution photocopies shrunken down to pocket size, slipping them into the 8 pockets, and viewing them through a magnifying glass ($1/S$) at destination.
\end{itemize}

\section{Core Mathematical Equations and Definitions}

\subsection{FP8 Quantization and Dequantization Map}
{\boldmath
\begin{equation}
X_{\text{quant}} = \operatorname{clamp}\left( \lfloor S \cdot X \rceil, \, -X_{\max}, \, X_{\max} \right), \quad X_{\text{recon}} = \frac{1}{S} X_{\text{quant}}
\end{equation}
}
\begin{itemize}
    \item \textbf{Formal Definition}: Symmetric uniform affine discretization and inverse restoration mapping.
    \item \textbf{What It Means}: Discretizes float representations into 8-bit bins and restores them via inverse scalar scale.
    \item \textbf{Why It Is Important}: Core mathematical operator executed by Transformer Engine kernels.
\end{itemize}

\subsection{Delayed Scale Update with Rolling History}
{\boldmath
\begin{equation}
S_{t+1} = \frac{X_{\max}^{\text{FP8}}}{\left( \max_{i=0}^{H-1} a_{\max}^{(t-i)} \right) \cdot 2^M}
\end{equation}
}
\begin{itemize}
    \item \textbf{Formal Definition}: Scale calculation based on historical peak tensor amplitudes.
    \item \textbf{What It Means}: Prepares next step's scale factor asynchronously during current execution.
    \item \textbf{Why It Is Important}: Avoids pipeline stalls and device synchronizations.
\end{itemize}

\section{Rigorous Mathematical Analysis Checklist}

\begin{itemize}
    \item \textbf{Notation}: $X \in \mathbb{R}^N$, $S$ (scaling factor), $X_{\max}^{\text{E4M3}} = 448$, $X_{\max}^{\text{E5M2}} = 57344$, $H$ (history window).
    \item \textbf{Epistemic Status}: Micikevicius et al. (NeurIPS 2022). OCP standard deployed across NVIDIA Hopper / Blackwell.
    \item \textbf{Dependency Graph}: Tensor $X \to \text{Evaluate } a_{\max} \to \text{Quantize with } S_t \to \text{GEMM in FP8} \to \text{Update } S_{t+1}$.
    \item \textbf{Dimensions}: Tensor length $N$, history length $H \in [8, 32]$.
    \item \textbf{Regimes}: Extreme-scale pre-training of frontier 70B+ LLMs.
    \item \textbf{Invariants}: Quantized representations never exceed format bounds $[-X_{\max}, X_{\max}]$.
    \item \textbf{Isomorphisms}: Isomorphic to dynamic floating-point block floating-point (BFP) quantization.
\end{itemize}

\section{Theoretical Theorems and Formal Proofs}

\begin{theorem}[Quantization Error Variance Bound of FP8 Formats]
Let input entries $x$ be independent zero-mean random variables with variance $\sigma^2$ scaled by factor $S$. Assuming uniform quantization noise $\epsilon \in [-0.5, 0.5]$ on the scaled integer grid, the mean squared quantization error is bounded by:
{\boldmath
\begin{equation}
\mathbb{E}\left[ (x - x_{\text{recon}})^2 \right] \le \frac{1}{12 S^2} + \mathcal{P}_{\text{clip}}(S)
\end{equation}
}
where $\mathcal{P}_{\text{clip}}(S) = \mathbb{E}\left[ (|x| - X_{\max}/S)^2 \cdot \mathbb{I}(|x| > X_{\max}/S) \right]$ is the clipping distortion.
\end{theorem}

\begin{proof}
Let $y = S \cdot x$. For non-clipped values $|y| \le X_{\max}$, rounding to the nearest integer introduces quantization noise $e = y - \lfloor y \rceil$.
By standard quantization theory (Widrow's theorem), for sufficiently fine quantization grids, $e$ is approximately uniformly distributed on $[-0.5, 0.5]$:
{\boldmath
\begin{equation}
\mathbb{E}[e] = 0, \quad \mathbb{E}[e^2] = \int_{-0.5}^{0.5} u^2 du = \left[ \frac{u^3}{3} \right]_{-0.5}^{0.5} = \frac{1}{12}
\end{equation}
}
Dequantizing yields error $\Delta x = \frac{e}{S}$:
{\boldmath
\begin{equation}
\mathbb{E}[(\Delta x)^2 \mid \text{no clip}] = \frac{1}{S^2} \mathbb{E}[e^2] = \frac{1}{12 S^2}
\end{equation}
}
For values exceeding $X_{\max}/S$, saturation clipping occurs:
{\boldmath
\begin{equation}
\Delta x_{\text{clip}} = x - \frac{X_{\max}}{S} \operatorname{sgn}(x)
\end{equation}
}
Combining both terms by the law of total expectation:
{\boldmath
\begin{equation}
\mathbb{E}\left[ (x - x_{\text{recon}})^2 \right] \le \frac{1}{12 S^2} + \mathbb{E}\left[ \left( |x| - \frac{X_{\max}}{S} \right)^2 \mathbb{I}\left(|x| > \frac{X_{\max}}{S}\right) \right]
\end{equation}
}
This completes the proof, showing that choosing scale $S$ balances quantization rounding noise $\mathcal{O}(1/S^2)$ against clipping saturation error.
\end{proof}

\begin{theorem}[Sufficient Condition for Non-Overflow Under Delayed Scaling]
If the sequence of tensor maximum values $a_{\max}^{(t)}$ has bounded step-to-step growth ratio $\frac{a_{\max}^{(t)}}{a_{\max}^{(t-1)}} \le 1 + \rho$, and history window $H \ge 1$, then choosing headroom margin $M \ge \log_2(1 + \rho)$ guarantees zero clipping overflow at step $t$:
{\boldmath
\begin{equation}
S_t \cdot a_{\max}^{(t)} \le X_{\max}^{\text{FP8}}
\end{equation}
}
\end{theorem}

\begin{proof}
By definition of the delayed scaling rule:
{\boldmath
\begin{equation}
S_t = \frac{X_{\max}^{\text{FP8}}}{\left( \max_{k \in [1, H]} a_{\max}^{(t-k)} \right) \cdot 2^M} \le \frac{X_{\max}^{\text{FP8}}}{a_{\max}^{(t-1)} \cdot 2^M}
\end{equation}
}
Multiplying by $a_{\max}^{(t)}$:
{\boldmath
\begin{equation}
S_t \cdot a_{\max}^{(t)} \le \frac{X_{\max}^{\text{FP8}}}{a_{\max}^{(t-1)} \cdot 2^M} \cdot a_{\max}^{(t)} = X_{\max}^{\text{FP8}} \cdot \frac{a_{\max}^{(t)}}{a_{\max}^{(t-1)}} \cdot 2^{-M}
\end{equation}
}
Given that $\frac{a_{\max}^{(t)}}{a_{\max}^{(t-1)}} \le 1 + \rho$ and setting $M \ge \log_2(1 + \rho) \implies 2^{-M} \le \frac{1}{1 + \rho}$:
{\boldmath
\begin{equation}
S_t \cdot a_{\max}^{(t)} \le X_{\max}^{\text{FP8}} \cdot (1 + \rho) \cdot \frac{1}{1 + \rho} = X_{\max}^{\text{FP8}}
\end{equation}
}
Thus, the scaled peak value does not exceed the format maximum, guaranteeing zero overflow.
\end{proof}

\section{Foundational Architectural Questions and Epistemic Meta-Questions}

\subsection{Architectural Question 1: Per-Tensor vs. Per-Block Scaling}
\textbf{Question:} Why is per-block FP8 scaling (fine-grained scaling) increasingly favored over per-tensor scaling?\\
\textbf{Answer:} Large models exhibit activation outliers (isolated channels with magnitudes $20\times\text{--}50\times$ larger than normal). In per-tensor scaling, a single outlier inflates $a_{\max}$, forcing scale $S$ down and squashing all normal activations into sub-bit noise. Per-block scaling applies independent scales per $1 \times 16$ or $1 \times 32$ tile, isolating outliers to their local blocks without degrading the rest of the tensor.

\subsection{Architectural Question 2: Memory vs. Compute Gains in FP8}
\textbf{Question:} Does FP8 training reduce optimizer memory?\\
\textbf{Answer:} No. FP8 speeds up matrix multiplications on tensor cores and halves activation cache memory, but master weights, gradients, and Adam moments ($m, v$) remain in FP32/BF16. To reduce optimizer memory, ZeRO (#234) or 8-bit Adam must be applied.

\section{Algorithmic Implementation Specification}

\begin{algorithm}[H]
\caption{FP8 Delayed Quantization and Dynamic Scale Update}
\begin{algorithmic}[1]
\Require Tensor $\mathbf{X} \in \mathbb{R}^N$, format $\in \{\text{e4m3}, \text{e5m2}\}$, scale $S$, history $\mathbf{H}$, mode, margin $M$
\Ensure Quantized vector $\mathbf{Q} \in \mathbb{Z}^N$, dequantized vector $\hat{\mathbf{X}} \in \mathbb{R}^N$, next scale $S_{\text{next}}$, SNR
\State $X_{\max} \gets (\text{format} == \text{"e4m3"}) \; ? \; 448.0 : 57344.0$
\State $a_{\max} \gets \max_{i} |\mathbf{X}[i]|$
\State $S_{\text{act}} \gets (\text{mode} == \text{"current"}) \; ? \; X_{\max} / (a_{\max} \cdot 2^M) : S$
\State Allocate $\mathbf{Q} \in \mathbb{Z}^N, \quad \hat{\mathbf{X}} \in \mathbb{R}^N$
\For{$i = 0$ \textbf{to} $N-1$}
    \State $v \gets \mathbf{X}[i] \cdot S_{\text{act}}$
    \State $q \gets \operatorname{round}(\max(-X_{\max}, \min(X_{\max}, v)))$
    \State $\mathbf{Q}[i] \gets q, \quad \hat{\mathbf{X}}[i] \gets q / S_{\text{act}}$
\EndFor
\State Update history $\mathbf{H}' \gets \mathbf{H} \cup [a_{\max}]$ (retaining last 16 values)
\State $S_{\text{next}} \gets X_{\max} / (\max(\mathbf{H}') \cdot 2^M)$
\State $\text{SNR} \gets 10 \log_{10}\left( \sum \mathbf{X}[i]^2 / \max(10^{-12}, \sum (\mathbf{X}[i] - \hat{\mathbf{X}}[i])^2) \right)$
\State \Return $(\mathbf{Q}, \hat{\mathbf{X}}, S_{\text{next}}, \text{SNR})$
\end{algorithmic}
\end{algorithm}

\section{Big-$\Theta$ Complexity Analysis}

\begin{itemize}
    \item \textbf{Quantization Pass Time}: $\Theta(N)$ arithmetic clamp and rounding operations.
    \item \textbf{Dequantization Pass Time}: $\Theta(N)$ multiplication operations.
    \item \textbf{History Update Time}: $\Theta(|H|)$ constant operations ($|H| \le 16$).
    \item \textbf{Total Time Complexity}: $\Theta(N)$ operations per tensor.
    \item \textbf{Memory Storage Space}: $\Theta(N)$ 8-bit integers (saving 50\% memory over FP16).
\end{itemize}

\section{Single Canonical Repository Reference}

The complete deterministic scalar implementation, verification suite, and unit test benchmarks are published at:
\begin{center}
\url{https://github.com/Chief-Strategist-J/llm-obs-infra}
\end{center}

\end{document}
